\manualchapter{Wave editing and patches}{sec:operation}
\subsection{Draw, blend and save}
Each of the five editors stores 32 samples with 16 possible levels. The bank
is shared across all polyphonic channels. Drawing supports undo/redo; patch
and preset saves include the samples. Reset restores the supplied sine,
pulse, rising ramp, distorted triangle and falling ramp; randomize replaces
them with smoothed random waves. Chip phase is not saved.

For a changing lead, draw a smooth wave in slot 1 and a jagged wave in slot 2.
Set Wave Morph to 1 and apply a slowly rising CV from 0 to 1.4\,V. Tune the
pulse voices an octave below and lower their levels to add a stable base.
For percussion, mute the pitched voices, send an envelope to Noise Level,
and sequence Noise Period in 2\,V increments.

\subsection{Output routing and polyphony}
A voice adds the preceding output only when that output is unpatched. Each
stage clips to approximately $\pm5$\,V; the final output is not a headroom-free
mixer. Patch an earlier output separately to remove its contribution from
the following mix, or reduce source levels.

The widest input selects 1--16 independent chip instances. Inputs generally
read the matching polyphonic channel rather than broadcasting mono CV.
Match pitch and level channel counts when playing chords. All instances use
the same five drawn waves and panel settings.

If the Wave output is silent, check its Level and CV: setting 0 is off. If
morph seems inactive near zero, move it above 1. If low-frequency modulation
sounds stepped, the wave volume has only four levels and the pulse/noise
volumes have sixteen. Use an external VCA where that stepping is unwanted.
